\documentclass[conference]{IEEEtran}
\IEEEoverridecommandlockouts

\usepackage[T1]{fontenc}
\usepackage[utf8]{inputenc}
\usepackage{graphicx}
\usepackage{booktabs}
\usepackage{amsmath}
\usepackage{array}
\usepackage{makecell}
\usepackage{multirow}
\usepackage{float}
\usepackage{cite}
\usepackage[final]{microtype}
\usepackage{url}
\usepackage[hidelinks]{hyperref}
\setlength{\emergencystretch}{3em}

% Slightly tighten table row padding to fit the page limit
% (purely cosmetic; no table content, columns, or margins are changed)
\renewcommand{\arraystretch}{0.95}

% Tighten bibliography vertical spacing to fit the page limit
% (purely cosmetic; no reference content is changed)
\let\oldthebibliography\thebibliography
\renewcommand{\thebibliography}[1]{%
  \oldthebibliography{#1}%
  \setlength{\itemsep}{0pt}%
  \setlength{\parskip}{0pt}%
  \setlength{\parsep}{0pt}%
}

\begin{document}

\title{In-Context Learning: Chain-of-Thought vs.\ Tree-of-Thoughts
Prompting on Multi-Step Reasoning}

\author{
\IEEEauthorblockN{O\u{g}uzhan \c{C}akal}
\IEEEauthorblockA{\textit{Department of Computer Engineering} \\
\textit{\.{I}zmir Institute of Technology (\.{I}YTE)}\\
\.{I}zmir, Turkey \\
oguzhancakal@std.iyte.edu.tr}
\and
\IEEEauthorblockN{Sercan Utku Temelci}
\IEEEauthorblockA{\textit{Department of Computer Engineering} \\
\textit{\.{I}zmir Institute of Technology (\.{I}YTE)}\\
\.{I}zmir, Turkey \\
sercanutkutemelci@std.iyte.edu.tr}
}

\maketitle

\begin{abstract}
We present an empirical comparison of four In-Context Learning (ICL)
prompting strategies on the GSM8K grade-school mathematics
benchmark---Zero-Shot Chain-of-Thought (CoT), Few-Shot CoT (8-shot),
Tree-of-Thoughts (ToT) with BFS at two search depths, and
Self-Consistency voting---using Llama~3.1~8B via the Groq API without
fine-tuning.
Zero-Shot CoT achieves \textbf{95.0\%}, Self-Consistency
\textbf{92.0\%}, and ToT (BFS-2, depth=6) and Few-Shot CoT each
\textbf{85.0\%} on the same 100-problem seed-42 sample.
A single answer-extraction format instruction raises Zero-Shot CoT by
$7$~pp (88.0\%~$\rightarrow$~95.0\%) and reorders the strategies---a
parsing confound, not a reasoning gain---placing it alone on the
accuracy--cost Pareto frontier at the 8B scale.
A depth ablation yields only $+4.0$~pp at $2.1\times$ cost, indicating
diminishing returns when the evaluator is the same 8B backbone.
A cross-model replication on Claude Haiku~4.5 with full evaluation
identity preserved collapses the inter-strategy gap from $14$ to
$2$~pp, revealing a \emph{ceiling effect}: deliberation- and
ensemble-style strategies pay off only when the base model has headroom
to lose and identifiable wrong branches to prune.
A pairwise Cohen's $\kappa$ analysis confirms this---ToT depths on
Llama reach near-zero agreement ($\kappa{=}0.011$) versus perfect
agreement on Haiku ($\kappa{=}1.000$).
\end{abstract}

\begin{IEEEkeywords}
Chain-of-thought, GSM8K, in-context learning, large language models,
multi-step reasoning, self-consistency, tree-of-thoughts
\end{IEEEkeywords}

%------------------------------------------------------------------------
\section{Introduction}
\label{sec:intro}

Large Language Models (LLMs) have demonstrated remarkable capabilities
in natural language tasks through In-Context Learning (ICL), where the
model adapts to new tasks from demonstrations provided in the prompt
without updating model weights~\cite{brown2020gpt3}.
This ability to generalise from a handful of examples has made ICL
the dominant paradigm for practical LLM deployment, motivating active
research into how the structure and content of prompts influence
multi-step reasoning quality.
Among ICL techniques, Chain-of-Thought (CoT) prompting has emerged as
a particularly effective approach, guiding models to produce
intermediate reasoning steps before arriving at a final
answer~\cite{wei2022cot}.
The zero-shot variant~\cite{kojima2022zeroshot} further simplified
this by showing that a single trigger phrase suffices for
instruction-tuned models to produce step-by-step reasoning, suggesting
that structured reasoning has been internalised during pre-training.

Tree-of-Thoughts (ToT) extends this paradigm by exploring multiple
reasoning branches simultaneously and pruning low-quality paths through
a separate evaluation step~\cite{yao2023tot}.
While ToT has shown impressive gains over CoT on GPT-4, its
effectiveness on smaller, open-weight models remains unexplored.
A central open question is whether the multi-branch search ToT
provides translates to accuracy gains when the evaluator model is the
same small backbone---or whether evaluator weakness undermines the
pruning mechanism entirely.

This paper presents a controlled empirical comparison of four
prompting strategies (Zero-Shot CoT, Few-Shot CoT (8-shot), ToT
with BFS, and Self-Consistency) on the GSM8K
benchmark~\cite{cobbe2021gsm8k}, using Llama~3.1~8B as a common
backbone, complemented by a ToT depth ablation and a cross-model
replication on Claude Haiku~4.5 (a frontier-tier model) under exact
sample, prompt, and evaluation identity.
Our three findings are summarised next as contributions.

\paragraph{Contributions}
This paper makes three empirical contributions, all derived from a
single controlled comparison rather than from new model training or
larger experiments:
\begin{itemize}
\item \textbf{A scale-dependent ceiling effect.} Strategy ordering is
not an intrinsic property of the prompting method but a function of
backbone capability: the inter-strategy accuracy spread collapses from
$14$~pp on Llama~3.1~8B to $2$~pp on Claude Haiku~4.5, and the
\emph{simplest} strategy on the strong backbone (Zero-Shot CoT, 97.0\%)
outperforms the \emph{most elaborate} strategy on the weak one
(Self-Consistency, 92.0\%). A pairwise Cohen's $\kappa$ analysis
confirms this independently
(Sections~\ref{sec:results:crossmodel}--\ref{sec:results:kappa}).
\item \textbf{Search depth is not the ToT bottleneck; evaluator
discrimination is.} Doubling ToT search depth (3$\rightarrow$6) buys
only $+4.0$~pp at $2.1\times$ inference cost on Llama and
$+0.0$~pp on Haiku, isolating the weak \emph{sure/maybe/impossible}
self-evaluator---not insufficient depth---as the factor that prevents
deliberate search from helping at the 8B scale
(Section~\ref{sec:results:depth}).
\item \textbf{Answer-extraction format is a confound in CoT
benchmarking.} A single output-format instruction raises Zero-Shot CoT
by $+7$~pp (88.0\%~$\rightarrow$~95.0\%) without changing the reasoning
process, showing that accuracy gaps attributed to a prompting
\emph{method} can partly reflect answer-extraction \emph{format} and
must be controlled when comparing strategies
(Section~\ref{sec:setup}).
\end{itemize}

%------------------------------------------------------------------------
\section{Related Work}
\label{sec:related}

The foundational prompting strategy for multi-step reasoning is
Chain-of-Thought (CoT), introduced by Wei et al.~\cite{wei2022cot},
who demonstrated that providing few-shot examples with step-by-step
solutions substantially improves LLM performance on arithmetic and
symbolic reasoning.
Kojima et al.~\cite{kojima2022zeroshot} subsequently showed that
simply appending \textit{``Let's think step by step''} to a query
elicits CoT reasoning in a zero-shot setting, establishing a
lightweight baseline that avoids hand-crafted demonstrations and
anchors our Zero-Shot CoT configuration.

Building directly on CoT, Yao et al.~\cite{yao2023tot} proposed
Tree-of-Thoughts (ToT), which frames reasoning as a search problem
over a tree of intermediate ``thoughts''.
At each node, the model generates multiple candidate continuations,
evaluates them with a separate scoring prompt (assigning
\textit{sure/maybe/impossible}), and retains only the most promising
branches using BFS or DFS.
Their experiments with GPT-4 showed substantial gains on tasks
requiring deliberate planning; however, all results used models with
100B+ parameters, leaving the effectiveness of the same mechanism on
smaller open-weight evaluators as an open question.

A complementary ensemble approach was proposed by Wang et
al.~\cite{wang2022selfconsistency}, who showed that sampling diverse
CoT paths at temperature $>0$ and selecting the final answer by
majority vote improves CoT accuracy by 10--20~pp without changing the
prompt.
Self-Consistency thus occupies a middle ground between the
minimal-cost single-chain CoT and the expensive multi-branch ToT.

All three strategies are instances of the In-Context Learning paradigm
of Brown et al.~\cite{brown2020gpt3}---generalising to unseen tasks from
a few in-context demonstrations without gradient updates---which makes
them directly comparable under a fixed model and prompt template.
Other structured-prompting variants (Automatic
CoT~\cite{zhang2022autocot}, which clusters questions to auto-generate
demonstrations, and Least-to-Most~\cite{zhou2022ltm}, which decomposes a
problem into ordered subproblems) map out a broader design space; we
focus on the four most widely adopted strategies for a controlled
comparison.

%------------------------------------------------------------------------
\section{Methodology}
\label{sec:method}

\subsection{Zero-Shot Chain-of-Thought}
The problem is presented with the trigger phrase
\textit{``Let's think step by step.''}~\cite{kojima2022zeroshot}
plus the output-format instruction
\textit{``End your response with: The answer is X.''}
to prevent trailing-modifier parse errors
(see Section~\ref{sec:setup}, \emph{Answer-extraction format fix}).

\subsection{Few-Shot Chain-of-Thought (8-shot)}
Eight worked examples from Wei et al.~\cite{wei2022cot} are prepended
to the query, each containing a full reasoning chain ending with
\textit{``The answer is X.''}\ followed by the test problem.

\subsection{Tree-of-Thoughts (BFS, branching=2, depth=6)}
\label{sec:method:totfixed}
ToT (\texttt{src/run\_tot\_fixed.py}) explores a tree of intermediate
``thoughts''.
At each depth level, two candidate next-step thoughts are generated at
temperature~$0.7$; a separate evaluator prompt scores each as
\textit{sure}, \textit{maybe}, or \textit{impossible}, and only the
former two are expanded.
After BFS completion, a dedicated step extracts the final answer from
the best leaf node.
The configuration (branching factor 2, depth 6) uses $\sim$45 API calls
per problem.

\subsection{Self-Consistency (5 paths)}
Few-Shot CoT is sampled 5 times per problem with temperature~$= 0.7$.
The final answer is selected by majority vote over the 5 sampled
numeric answers~\cite{wang2022selfconsistency}.

\subsection{Cross-Model Evaluation Protocol}
\label{sec:method:crossmodel}
To test whether strategy orderings generalise beyond the 8B scale, we
re-run all five configurations on Claude Haiku~4.5 with identical
100-problem sample, prompts, parser, and gold-label correction, so any
accuracy difference is attributable to the backbone alone.
The three single-call strategies are batched through the Anthropic
\texttt{messages.batches} API (50\% discount); ToT runs synchronously
because its BFS state cannot be flattened into batch items.

%------------------------------------------------------------------------
\section{Experimental Setup}
\label{sec:setup}

\subsection{Dataset}
We use GSM8K~\cite{cobbe2021gsm8k}, a benchmark of 8,792 grade-school
math word problems requiring 2--8 chained reasoning steps.
All five configurations (Zero-Shot CoT, Few-Shot CoT, ToT depth=3,
ToT depth=6, and Self-Consistency) use the same 100-problem random
sample of GSM8K's 1,319-item test split (seed~$= 42$).
One item in our sample contains a gold-label annotation error
(see the \emph{GSM8K annotation error} paragraph below); we manually
verified the correct value and use the corrected label in all
evaluations, preserving the original sample sizes.

\subsection{Model}
All experiments use Llama~3.1~8B (\texttt{llama-3.1-8b-instant}) via
the Groq inference API.
No fine-tuning is applied; the model is used in its base instruction
configuration.

\subsection{Evaluation}
The primary metric is \textit{exact-match accuracy}: the predicted
final number is compared against the gold answer after stripping commas
and converting to float.
Answers are extracted by a regular expression capturing the gold-format
marker \texttt{\#\#\#\#}, falling back to the last number otherwise.
We also report the \textit{parse failure rate} (responses yielding no
numeric prediction).
Llama~3.1~8B never spontaneously produced the \texttt{\#\#\#\#} marker,
so every response used the last-number fallback, which succeeded on all
500 evaluated responses (parse failure rate~$= 0\%$).

\subsection{Answer-Extraction Format Fix}
An initial Zero-Shot CoT run sometimes ended with a trailing-modifier
number rather than the substantive answer (e.g.\ ``\textit{36 hours
in 4 weeks}'' $\rightarrow$ parser captures \texttt{4} instead of \texttt{36}).
Appending ``\textit{End your response with: The answer is X.}'' to the
prompt raised Zero-Shot CoT from \textbf{88.0\%} to \textbf{95.0\%}
($+7$\,pp), concentrated in the hard bucket.
All Zero-Shot CoT numbers use this format-aligned prompt; the other
strategies already produce unambiguous final answers via their
demonstrations or dedicated extraction prompts.
Crucially, without this format alignment Zero-Shot CoT (88.0\%) would
appear merely tied with Few-Shot CoT and ToT (both 85.0\%): the
strategy ordering itself is altered by an extraction-format change that
leaves the reasoning process untouched.
Reported accuracy gaps between prompting strategies can therefore
partly reflect answer-extraction format rather than reasoning ability,
and the format must be held fixed---or reported---for a fair comparison
(see Section~\ref{sec:discussion}).

\subsection{GSM8K Annotation Error (Gold-Label Correction)}
Manual inspection identified one annotation error in our sample (the
\textit{carnival} item): GSM8K's own calculator-annotated solution,
\texttt{750+430+400+700=2280}, substitutes Maryam's \$400 excess over
Sarah for Sarah's own \$300 contribution; the correct total is
$430+750+300+700=\$2{,}180$. We adopt the corrected label for all
reported accuracies. On Llama, Zero-Shot CoT and ToT depth=6 shift
from incorrect to correct; on Haiku, all five strategies gain a
uniform $+1$\,pp. In both cases the correction leaves strategy ranking
unchanged, consistent with the $\sim$1--2\% annotation noise floor of
GSM8K~\cite{cobbe2021gsm8k}.

\subsection{Infrastructure}
Rate-limit constraints on the Groq free tier required an 8-second
inter-call pause (CoT) and a 30-second inter-problem pause (ToT).
All random seeds are fixed at 42 for reproducibility.
All code and result files are publicly available.\footnote{\url{https://github.com/OguzhanCKL/CENG467_In-Context_Learning}}

\subsection{Secondary Model and Cross-Model Setup}
\label{sec:setup:crossmodel}
The secondary backbone is Claude Haiku~4.5
(\texttt{claude-haiku-4-5-20251001}) via the Anthropic Python SDK,
run on the same seed-42 sample, prompts, parser, and carnival
gold-label correction as Llama, with batching as described in
Section~\ref{sec:method:crossmodel} (700 batched requests).
Realized total cost: \$3.26 across all five strategies
(\$0.52 batch + \$0.72 ToT-d3 + \$2.02 ToT-d6).

%------------------------------------------------------------------------
\section{Results}
\label{sec:results}

\subsection{Main Results}

Table~\ref{tab:main} summarises the accuracy of the four strategies.

\begin{table*}[t]
\centering
\caption{Main results on GSM8K. Llama~3.1~8B, Groq API.
Accuracies use the GSM8K gold labels with one item corrected (see
Section~\ref{sec:setup}, \emph{GSM8K annotation error}).
$\bar{T}$ denotes the average response length in tokens (one token
$\approx$ four characters), aggregated across all API calls per problem;
TPC = tokens per correct answer (lower is better).}
\label{tab:main}
\begin{tabular}{@{}lcccccl@{}}
\toprule
\textbf{Strategy} & \textbf{Accuracy} & \textbf{N} &
\textbf{API Calls} & \boldmath$\bar{T}$ & \textbf{TPC} & \textbf{Status} \\
\midrule
Zero-Shot CoT                     & \textbf{95.0\%} & 100 & $\sim$1  & 189  & 199  & Complete \\
Few-Shot CoT (8-shot)             & \textbf{85.0\%} & 100 & $\sim$1  & 145  & 171  & Complete \\
Tree-of-Thoughts (BFS-2, depth=3) & \textbf{81.0\%} & 100 & $\sim$21 & 315  & 389  & Complete (ablation) \\
Tree-of-Thoughts (BFS-2, depth=6) & \textbf{85.0\%} & 100 & $\sim$45 & 561  & 660  & Complete \\
Self-Consistency (5 paths)        & \textbf{92.0\%} & 100 & 5        & 738  & 802  & Complete \\
\bottomrule
\end{tabular}
\end{table*}

Beyond accuracy, we report two efficiency metrics that have become
standard in recent reasoning benchmarks~\cite{ockbench2025,tale2024}:
the average per-problem token budget $\bar{T}$ (a model-agnostic proxy
for compute cost) and \textit{tokens per correct answer}
($\text{TPC} = \bar{T} / \text{accuracy}$), which captures the
practical cost of \emph{useful} reasoning---a method achieving high
accuracy but spending disproportionately many tokens may be
impractical~\cite{yao2023tot,ockbench2025}.
Few-Shot CoT is the most token-efficient strategy (171 TPC), followed
closely by Zero-Shot CoT (199); both ToT configurations spend
2--4$\times$ more tokens per correct answer (389 and 660), and
Self-Consistency---despite the highest accuracy---is the most expensive
at 802~TPC, consistent with its 5$\times$ sample budget.

Zero-Shot CoT achieves 95.0\% (after the answer-extraction
format fix, Section~\ref{sec:setup}) and substantially
outperforms Few-Shot CoT (85.0\%), reversing the typical
expectation that more examples improve performance.
This suggests that Llama~3.1~8B has already internalised strong
step-by-step reasoning during pre-training, and the 8-shot
demonstrations introduce noise rather than useful signal for this
sample.

ToT (depth=6) matches Few-Shot CoT at 85.0\% but trails Zero-Shot CoT
by $10.0$\,pp; Section~\ref{sec:results:crossmodel} tests whether these
orderings generalise to a stronger backbone.

\subsection{Depth Ablation for Tree-of-Thoughts}
\label{sec:results:depth}

To isolate the effect of search depth, we ran ToT again on the same
100-problem sample with \texttt{depth=3} (branching factor unchanged at
2), mirroring the short-horizon depth of the original ToT
paper~\cite{yao2023tot} against our deeper \texttt{depth=6} on identical
problems and the same evaluator prompt.
Table~\ref{tab:depth} reports the result and the accuracy--cost
trade-off.

\begin{table}[t]
\centering
\caption{ToT depth ablation on GSM8K (Llama~3.1~8B, BFS-2,
identical evaluator prompt). Accuracies use the corrected GSM8K labels
(see Section~\ref{sec:setup}).}
\label{tab:depth}
\begin{tabular}{@{}lcccc@{}}
\toprule
\textbf{Configuration} & \textbf{Accuracy} & \textbf{N} &
\makecell{\textbf{API calls}\\\textbf{per problem}} &
\makecell{\textbf{Wall-clock}\\\textbf{per problem}} \\
\midrule
ToT BFS-2, depth=3 & 81.0\% & 100 & $\sim$21 & $\sim$3 min \\
ToT BFS-2, depth=6 & 85.0\% & 100 & $\sim$45 & $\sim$6 min \\
\midrule
$\Delta$ (depth=6 $-$ depth=3) & \textbf{+4.0\,pp} & -- &
\textbf{+114\%} & \textbf{+100\%} \\
\bottomrule
\end{tabular}
\end{table}

The $+4.0$\,pp gain is small relative to the $2.1\times$ cost increase.
Unlike GPT-4's 100B+ evaluator~\cite{yao2023tot}, our 8B
\textit{sure/maybe/impossible} scorer is weakly discriminative---most
partial thoughts receive \textit{sure} regardless of correctness,
collapsing BFS toward a near-linear chain.
GSM8K problems (2--8 steps) also fit comfortably within depth=3,
so additional depth helps only a small tail of multi-step problems.
The bottleneck at the 8B scale is therefore evaluator discrimination
rather than search depth.

\subsection{Self-Consistency: Agreement and Pass@k}
\label{sec:results:passk}

We summarise two Self-Consistency diagnostics in text for brevity.
\emph{Vote agreement:} all five paths agreed unanimously on 65\% of
problems and reached a 4/5 or 3/5 majority on a further 30\%, so 95\%
of problems admit a clean decision (88.2\% average majority margin),
confirming that 5 paths suffice for reliable consensus.
\emph{Pass@k}~\cite{chen2021codex} (the fraction of problems where at
least one of the first $k$ paths is correct): Pass@1 $=86.4\%$ closely
tracks greedy Few-Shot CoT (85.0\%), so a single stochastic sample adds
no consistent gain, whereas Pass@5 reaches $99.0\%$.
This exposes a $+7.0$\,pp \textit{verification gap} over majority
voting (92.0\%): in 7 of 100 problems the correct answer is among the
five sampled paths but loses the vote, indicating that a smarter
aggregator---not more samples---is the main opportunity for further
gains.

\subsection{Accuracy Stratified by Problem Complexity}
\label{sec:results:stratified}

We adopt the Cobbe et al.~\cite{cobbe2021gsm8k} difficulty split:
problems are bucketed by gold-solution operation count into
\textit{easy} ($<3$), \textit{medium} ($=3$), and \textit{hard}
($>3$) groups (30/33/37 in our sample).
Table~\ref{tab:strat} reports per-bucket accuracy.

\begin{table*}[t]
\centering
\caption{Accuracy stratified by problem complexity
(Cobbe et al.~\cite{cobbe2021gsm8k} split, gold-corrected; N=100,
buckets 30/33/37) on both backbones. Llama columns are discussed in
this section; Haiku columns in Section~\ref{sec:results:crossmodel}.
All values are percentages; easy/medium/hard $=$ $<$3 / $=$3 / $>$3
gold operations.}
\label{tab:strat}
\setlength{\tabcolsep}{5pt}
\small
\begin{tabular}{@{}lcccc@{\hspace{1.6em}}cccc@{}}
\toprule
& \multicolumn{4}{c}{\textbf{Llama~3.1~8B}} &
  \multicolumn{4}{c}{\textbf{Haiku~4.5}} \\
\cmidrule(lr){2-5}\cmidrule(l){6-9}
\textbf{Strategy} & \textbf{Easy} & \textbf{Med} & \textbf{Hard} & \textbf{Overall}
& \textbf{Easy} & \textbf{Med} & \textbf{Hard} & \textbf{Overall} \\
\midrule
Zero-Shot CoT              & \textbf{96.7} & \textbf{93.9} & \textbf{94.6} & \textbf{95.0} & \textbf{100} & 97.0 & 94.6 & 97.0 \\
Few-Shot CoT (8-shot)      & 90.0 & 87.9 & 78.4 & 85.0 & \textbf{100} & \textbf{100} & 94.6 & 98.0 \\
ToT (BFS-2, depth=3)       & 86.7 & 75.8 & 81.1 & 81.0 & \textbf{100} & \textbf{100} & 94.6 & 98.0 \\
ToT (BFS-2, depth=6)       & 90.0 & 84.8 & 81.1 & 85.0 & \textbf{100} & \textbf{100} & 94.6 & 98.0 \\
Self-Consistency (5 paths) & 96.7 & 93.9 & 86.5 & 92.0 & \textbf{100} & \textbf{100} & \textbf{97.3} & \textbf{99.0} \\
\bottomrule
\end{tabular}
\end{table*}

Three results stand out.
First, every strategy shows the expected easy-to-hard accuracy decay.
Second, Few-Shot CoT loses the most on hard problems ($-16.2$\,pp vs.\
Zero-Shot CoT), since its short 2--3-step demonstrations do not
transfer well to longer reasoning chains.
Third, ToT (depth=3) collapses on medium problems (75.8\%) because
3-operation problems exactly exhaust its depth budget, leaving no slack
for BFS recovery after an initial error.
After the format fix, Zero-Shot CoT ties Self-Consistency on easy and
medium buckets and \emph{exceeds} it on hard ($94.6\%$ vs.\ $86.5\%$,
$+8.1$\,pp)---Zero-Shot CoT with disciplined answer extraction is the
strongest strategy at the 8B scale across all difficulty buckets.

\subsection{Solution Length Ratio: Reasoning Verbosity and Overthinking}
\label{sec:results:slr}

Token cost ($\bar{T}$, TPC) does not reveal whether compute is
well-spent.
We quantify reasoning verbosity via the \textbf{Solution Length Ratio}
$\text{SLR} = (\text{model reasoning steps}) / (\text{gold solution
steps})$, where model steps count \texttt{``= NUMBER''} patterns and
gold steps are the GSM8K \texttt{<<...>>} annotations
($\text{SLR}{>}1$ indicates overthinking~\cite{wu2025whenmoreisless};
average gold length is 3.3 steps).
Self-Consistency is the most compact strategy ($1.27\times$), even
below Few-Shot CoT ($1.44\times$) and Zero-Shot CoT ($1.60\times$),
because averaging across five paths filters out the longest chains.
ToT (depth=6), by contrast, inflates solution length to $3.30\times$
(depth=3: $2.72\times$) while only matching Few-Shot CoT in
accuracy---a clean instance of the \textit{overthinking
tax}~\cite{wu2025whenmoreisless}: BFS expansions add depth without
semantic gain because the 8B evaluator cannot prune unproductive
branches.

\subsection{Cross-Model Comparison: Llama~3.1~8B vs.\ Haiku~4.5}
\label{sec:results:crossmodel}

To test whether strategy orderings reflect intrinsic prompting
properties or 8B-scale limitations, we re-ran every strategy on
Claude Haiku~4.5 with full evaluation identity preserved
(Sections~\ref{sec:method:crossmodel}--\ref{sec:setup:crossmodel}).

\paragraph{Headline accuracy}
Table~\ref{tab:crossmodel} reports overall accuracy on the 100-problem
seed-42 sample under the gold-corrected labels for both models.
\begin{table*}[t]
\centering
\caption{Cross-model accuracy (gold-corrected) on the same
100-problem GSM8K sample. ``Cost'' is the realized USD cost for the
full 100-problem run on Haiku~4.5; ZS/FS/SC share a single
700-request batch with the $50\%$ batch discount.}
\label{tab:crossmodel}
\begin{tabular}{@{}lccrc@{}}
\toprule
\textbf{Strategy} & \textbf{Llama~3.1~8B} & \textbf{Haiku~4.5} &
\boldmath$\Delta$ & \textbf{Haiku cost} \\
\midrule
Zero-Shot CoT             & 95.0\% & \textbf{97.0\%} & $+2.0$ &
\multirow{3}{*}{\$0.52 (batch)} \\
Few-Shot CoT (8-shot)     & 85.0\% & \textbf{98.0\%} & $+13.0$ & \\
Self-Consistency (5 paths)& 92.0\% & \textbf{99.0\%} & $+7.0$ & \\
\midrule
ToT (BFS-2, depth=3)      & 81.0\% & \textbf{98.0\%} & $+17.0$ & \$0.72 (sync) \\
ToT (BFS-2, depth=6)      & 85.0\% & \textbf{98.0\%} & $+13.0$ & \$2.02 (sync) \\
\midrule
\multicolumn{4}{r}{\textbf{Total realized cost on Haiku~4.5:}} &
\textbf{\$3.26} \\
\bottomrule
\end{tabular}
\end{table*}

Every strategy improves on the stronger backbone, as expected; the more
interesting observation is that the \emph{gaps between strategies
collapse}: the five strategies span $14$~pp on Llama (81.0\% to 95.0\%)
but only $2$~pp on Haiku (97.0\% to 99.0\%).
The simplest strategy on Haiku (Zero-Shot CoT, 97.0\%) already
outperforms the most elaborate strategy on Llama (Self-Consistency,
92.0\%) by $5$~pp---at GSM8K's difficulty, raw backbone capability
dominates prompting-strategy choice once the backbone is strong enough.

\paragraph{Stratified view}
The Haiku columns of Table~\ref{tab:strat} report the same Cobbe
difficulty buckets alongside the Llama numbers.
Four of five strategies achieve a perfect 63/63 on Easy and Medium;
Haiku's entire residual error budget lives in the Hard bucket.
Only Self-Consistency recovers one additional Hard answer (36/37 vs.\
35/37), mirroring its 8B advantage---Hard accuracy is the sole
discriminating metric once the backbone nears ceiling.

\paragraph{Cross-model depth ablation: ceiling effect}
Table~\ref{tab:crossmodel:depth} shows the sharpest single contrast:
depth=3$\rightarrow$6 gains $+4.0$~pp on Llama but $0.0$~pp on Haiku at
$2.8\times$ the token cost.
\begin{table}[t]
\centering
\caption{Cross-model ToT depth ablation ($\Delta$ = gold-corrected accuracy).}
\label{tab:crossmodel:depth}
\begin{tabular}{@{}lccc@{}}
\toprule
\textbf{Backbone} &
\textbf{depth=3} & \textbf{depth=6} & \boldmath$\Delta$ \\
\midrule
Llama~3.1~8B   & 81.0\% & 85.0\% & $+4.0$\,pp \\
Haiku~4.5      & 98.0\% & 98.0\% & $\phantom{+}0.0$\,pp \\
\midrule
Haiku token cost &
\multicolumn{3}{r}{depth=6 is $3.1\times$ input and $2.4\times$ output of depth=3} \\
\bottomrule
\end{tabular}
\end{table}
On Haiku neither depth beats Zero-Shot CoT on Hard (both 94.6\%);
ToT's depth parameter buys accuracy only when the base model has
headroom to lose and the evaluator has branches to discriminate.

\paragraph{Self-Consistency verifier saturation}
On Llama the verifier gap is $+7.0$\,pp (Pass@5 = 99\%,
majority-vote = 92\%): the correct answer is present in the five
paths but discarded by majority voting.
On Haiku the gap is $0.0$\,pp (Pass@5 = majority-vote = 99\%): on
the one remaining error, \emph{none} of the five paths produces the
correct answer.
Path diversity has collapsed (98/100 problems show unanimous 5/5
voting), confirming that Self-Consistency's headroom is a model-scale
artefact, not an intrinsic strategy benefit.

\subsection{Strategy Agreement: Pairwise Cohen's $\kappa$ Across Backbones}
\label{sec:results:kappa}

Accuracy alone does not reveal whether two strategies are correct on
the \emph{same} problems.
We therefore compute all pairwise \textbf{Cohen's $\kappa$}~\cite{cohen1960kappa}
over the same 100-problem gold-corrected sample, treating each strategy
as a binary classifier; $\kappa=0$ is chance, $\kappa=1$ is perfect
agreement (Landis--Koch ranges~\cite{landis1977kappa}; see also Artstein
and Poesio~\cite{artstein2008agreement} on agreement coefficients for
language data).

\begin{table}[t]
\centering
\caption{Pairwise Cohen's $\kappa$ between the five strategies
(100-problem seed-42 sample, gold-corrected).
\textbf{Lower triangle:} Llama~3.1~8B; \textbf{upper triangle:}
Haiku~4.5.}
\label{tab:kappa}
\setlength{\tabcolsep}{4pt}
\begin{tabular}{@{}lccccc@{}}
\toprule
& \textbf{ZS} & \textbf{FS} & \textbf{SC} & \textbf{d=3} & \textbf{d=6} \\
\midrule
ZS-CoT  & --     & 0.385 & 0.492 & \textbf{0.795} & \textbf{0.795} \\
FS-CoT  & 0.243  & --    & 0.662 & 0.490 & 0.490 \\
SC (5)  & 0.262  & \textbf{0.563} & --    & 0.662 & 0.662 \\
ToT d=3 & 0.095  & 0.435 & 0.207 & --    & \textbf{1.000} \\
ToT d=6 & 0.027  & 0.216 & 0.078 & \textbf{0.011} & --    \\
\bottomrule
\end{tabular}
\end{table}

The two backbones invert.
On Llama (lower triangle), the two ToT depths reach near-zero pairwise
agreement ($\kappa=0.011$) despite sharing the same BFS-2 algorithm and
evaluator prompt---BFS exploration at the 8B scale is stochastic rather
than deliberate---while Zero-Shot CoT shows only fair-to-slight
agreement with every other strategy ($\kappa \in [0.027, 0.262]$), the
genuine complementarity that drives the $+7.0$\,pp Pass@5 headroom
(Section~\ref{sec:results:passk}).
On Haiku (upper triangle) the picture reverses: the two ToT depths
reach $\kappa=1.000$ (perfect problem-by-problem concordance) and both
agree with Zero-Shot CoT at $\kappa=0.795$, so deliberate-search
machinery collapses into single-chain behaviour.
The overall range collapses from $[0.011,\,0.563]$ on Llama to
$[0.385,\,1.000]$ on Haiku: \emph{strategy independence itself decays
with model scale}.

\subsection{Ablation Study}
Two secondary Few-Shot CoT ablations (N=50, seed=42) are summarised in
text.
\emph{Temperature:} greedy decoding (88.0\%) outperforms stochastic
sampling at $T{=}0.7$ (86.0\%) by $2$\,pp, a marginal but consistent
advantage for deterministic decoding on arithmetic tasks.
\emph{Prompt sensitivity:} accuracy across three 8-shot example sets
spans $6$\,pp (82.0\%--88.0\%)---Set-1 underperforms (82.0\%) while the
canonical Wei et al.\ examples and Set-2 both reach 88.0\%---indicating
moderate sensitivity to example choice.

%------------------------------------------------------------------------
\section{Discussion}
\label{sec:discussion}

\vspace{4.5pt}\noindent\textbf{Zero-shot CoT wins; Self-Consistency no longer dominates.}

The 10-pp gap between Zero-Shot and Few-Shot CoT (Table~\ref{tab:main})
indicates that Llama~3.1~8B already internalised step-by-step reasoning
during pre-training, so the 8-shot demonstrations inject systematic
noise rather than signal---a prompt-format mismatch that is actively
harmful, not merely neutral.
Once the answer-extraction fix lifts Zero-Shot CoT to 95.0\%,
Self-Consistency (at $5\times$ the cost) is strictly Pareto-dominated:
its per-path Pass@1 (86.4\%) sits below the greedy chain.

\vspace{4.5pt}\noindent\textbf{Accuracy--cost trade-off.}

On a log-scale cost axis (average API calls per problem), Zero-Shot CoT
($\sim$1 call, 95.0\%) lies alone on the Pareto frontier; every other
configuration is dominated---Self-Consistency costs $5\times$ for
$-3.0$\,pp, Few-Shot CoT matches the cost at $-10.0$\,pp, and ToT costs
$21$--$45\times$ for $-10$ to $-14$\,pp (see $\bar{T}$/TPC in
Table~\ref{tab:main}).
Where accuracy is paramount and cost secondary, Self-Consistency is the
practical choice; otherwise Zero-Shot CoT offers the best
accuracy-per-call ratio.

\vspace{4.5pt}\noindent\textbf{ToT and the ceiling effect.}

ToT (depth=6) matches but never beats the CoT baselines on either
backbone, so the absence of a ToT advantage is not an artefact of an
undersized evaluator.
Two factors explain this on Llama: GSM8K problems resolve in 3--5 steps
that CoT already handles, and the \textit{sure/maybe/impossible}
evaluator (the same 8B model) is too weak to discriminate, collapsing
BFS toward a near-linear chain (the depth ablation buys only
$+4.0$\,pp at $2.1\times$ cost).
The Haiku replication sharpens this into a \emph{ceiling effect}: the
inter-strategy spread collapses from $14$~pp to $2$~pp, the
depth=3$\rightarrow$6 gain falls to $0.0$\,pp, and the
Self-Consistency verifier gap shrinks from $+7.0$ to $0.0$\,pp because
no sampled path recovers the problems voting gets wrong.
The pairwise $\kappa$ analysis confirms this independently
(Table~\ref{tab:kappa}): Llama's ToT depths agree at only
$\kappa=0.011$ (stochastic exploration) versus $\kappa=1.000$ on Haiku
(deterministic ceiling).
Deliberation- and ensemble-style strategies thus buy accuracy only when
the base model has headroom to lose and identifiable wrong branches to
prune.

\vspace{4.5pt}\noindent\textbf{Answer extraction is a confound, not a detail.}

Before the format fix, Zero-Shot CoT (88.0\%) was statistically
indistinguishable from Few-Shot CoT and ToT (depth=6) at 85.0\%, and a
reader would reasonably conclude that the three strategies perform
alike at the 8B scale.
After appending a single output-format instruction---which changes
nothing about how the model reasons---Zero-Shot CoT becomes the clear
leader at 95.0\%, and the gain concentrates in the hard bucket where
trailing-modifier numbers are most common.
The methodological implication is that, especially for smaller models,
benchmark rankings can be sensitive to answer-extraction format rather
than to the reasoning method under study; comparisons across prompting
strategies should hold the extraction format fixed or report it
explicitly, lest a parsing artefact masquerade as a method effect.

\vspace{4.5pt}\noindent\textbf{Error analysis and limitations.}

Across all strategies and both backbones, the dominant failure mode is
arithmetic: the model builds the correct reasoning structure but
miscomputes one step, and the error propagates---on Haiku every
residual error is of this kind, confirming that arithmetic brittleness
persists even at frontier scale, only attenuated.
The main limitations are: (1)~exact-match accuracy ignores reasoning
quality; (2)~the cross-model comparison covers only two families---adding
Mistral, Qwen, or Gemma at the 7--8B scale would separate pre-training
data quality from raw scale; and (3)~GSM8K's short arithmetic horizon
means these findings should not be extrapolated to long-horizon
planning benchmarks where ToT reported larger gains~\cite{yao2023tot}.

%------------------------------------------------------------------------
\section{Conclusion}
\label{sec:conclusion}

We evaluated Zero-Shot CoT, Few-Shot CoT, ToT (at two search depths),
and Self-Consistency on GSM8K across Llama~3.1~8B and Claude Haiku~4.5.
On Llama, Zero-Shot CoT with disciplined answer extraction (95.0\%) is
both the most accurate and the most cost-efficient strategy---a strict
Pareto dominator---while ToT only matches Few-Shot CoT and its depth
ablation yields diminishing returns ($+4.0$\,pp at $2.1\times$ cost)
due to weak 8B-evaluator discrimination rather than insufficient search
depth.
We further show that this format alignment is not a cosmetic
detail but a confound in CoT benchmarking: the same $+7$\,pp it
contributes reorders the strategies, so reported method gaps can
partly reflect answer-extraction format and must be controlled.
Replicating all five strategies on the stronger backbone with full
evaluation identity preserved exposes a \emph{ceiling effect}: the
inter-strategy gap collapses from $14$ to $2$~pp, the ToT depth gain
vanishes, and the Self-Consistency verifier gap closes---deliberation-
and ensemble-style strategies buy real accuracy only when the base
model has headroom to lose and identifiable wrong branches to prune.

\paragraph{Future work}
Replacing the categorical \textit{sure/maybe/impossible} evaluator with
a learned scorer is the most direct route to better ToT on 8B-class
models, and the $7.0$\,pp Pass@5--vote verification gap on Llama
suggests that an external verifier---rather than larger ensembles---is
the main opportunity at fixed sample budget.
Extending the comparison to additional open-weight families (Mistral,
Qwen, Gemma at the 7--8B scale) would disentangle pre-training data
quality from raw scale and test whether the ceiling effect generalises.

%------------------------------------------------------------------------
\begingroup
\small
\begin{thebibliography}{00}

\bibitem{brown2020gpt3}
T. Brown \textit{et al.}, ``Language models are few-shot learners,'' in
\textit{Adv. Neural Inf. Process. Syst. (NeurIPS)}, vol. 33, 2020,
pp. 1877--1901.

\bibitem{wei2022cot}
J. Wei \textit{et al.}, ``Chain-of-thought prompting elicits reasoning
in large language models,'' in \textit{Adv. Neural Inf. Process. Syst.
(NeurIPS)}, vol. 35, 2022, pp. 24824--24837.

\bibitem{kojima2022zeroshot}
T. Kojima, S. S. Gu, M. Reid, Y. Matsuo, and Y. Iwasawa, ``Large
language models are zero-shot reasoners,'' in \textit{Adv. Neural Inf.
Process. Syst. (NeurIPS)}, vol. 35, 2022, pp. 22199--22213.

\bibitem{yao2023tot}
S. Yao \textit{et al.}, ``Tree of thoughts: Deliberate problem solving
with large language models,'' in \textit{Adv. Neural Inf. Process. Syst.
(NeurIPS)}, vol. 36, 2023.

\bibitem{cobbe2021gsm8k}
K. Cobbe \textit{et al.}, ``Training verifiers to solve math word
problems,'' arXiv:2110.14168, 2021.

\bibitem{wang2022selfconsistency}
X. Wang \textit{et al.}, ``Self-consistency improves chain of thought
reasoning in language models,'' in \textit{Proc. Int. Conf. Learn.
Represent. (ICLR)}, 2023.

\bibitem{zhang2022autocot}
Z. Zhang, A. Zhang, M. Li, and A. Smola, ``Automatic chain of thought
prompting in large language models,'' in \textit{Proc. Int. Conf. Learn.
Represent. (ICLR)}, 2023.

\bibitem{zhou2022ltm}
D. Zhou \textit{et al.}, ``Least-to-most prompting enables complex
reasoning in large language models,'' in \textit{Proc. Int. Conf. Learn.
Represent. (ICLR)}, 2023.

\bibitem{ockbench2025}
OckBench Authors, ``OckBench: Measuring the efficiency of LLM
reasoning,'' arXiv:2511.05722, 2025.

\bibitem{tale2024}
T. Han \textit{et al.}, ``Token-budget-aware LLM reasoning,''
arXiv:2412.18547, 2024.

\bibitem{chen2021codex}
M. Chen \textit{et al.}, ``Evaluating large language models trained on
code,'' arXiv:2107.03374, 2021.

\bibitem{wu2025whenmoreisless}
Y. Wu \textit{et al.}, ``When more is less: Understanding
chain-of-thought length in LLMs,'' arXiv:2502.07266, 2025.

\bibitem{cohen1960kappa}
J. Cohen, ``A coefficient of agreement for nominal scales,''
\textit{Educational and Psychological Measurement}, vol. 20, no. 1,
pp. 37--46, 1960.

\bibitem{landis1977kappa}
J. R. Landis and G. G. Koch, ``The measurement of observer agreement
for categorical data,'' \textit{Biometrics}, vol. 33, no. 1,
pp. 159--174, 1977.

\bibitem{artstein2008agreement}
R. Artstein and M. Poesio, ``Inter-coder agreement for computational
linguistics,'' \textit{Computational Linguistics}, vol. 34, no. 4,
pp. 555--596, 2008.

\end{thebibliography}
\endgroup

\end{document}
